\section{Conclusion}

Cervical-spine MRI segmentation lacked a fairness audit across multiple
demographic axes. We conducted one on the CSpineSeg dataset ($1{,}254$ exams),
evaluating performance across sex, race, and age with anatomically informed
metrics.

The model is fair by every measure we applied: $0/63$ statistical tests reach
significance after FDR correction, all disparate impact ratios exceed $0.96$, and
demographics collectively explain $1.3\%$ of Dice variance. Replacing expert
reference labels with machine-generated ones flips the verdict for age --
not by inflating the gap but by collapsing within-group variance ${\sim}4\times$,
a false-confidence mode complementary to the false-magnitude effect of Parikh et
al.~\cite{parikh2025labelbias}. Training on silver labels does not amplify
demographic bias in this anatomy, unlike the amplification observed in breast
MRI~\cite{parikh2025labelbias}. Encoder probing confirms that the residual age
gradient reflects intrinsic anatomical difficulty -- older spines are harder to
segment -- rather than algorithmic harm, and that strong encoding of a protected
attribute (sex, AUROC $0.957$) need not produce any performance disparity.

These results carry two practical implications. First, reference-label provenance
is a first-order confounder in fairness evaluation: the same predictions, scored
against different labels, yield opposite conclusions. Any fairness claim must
state where its ruler comes from. Second, debiasing a model's learned
representations without first confirming an output-level disparity addresses a
problem that may not exist -- our model strongly encodes sex yet exhibits no
sex-based performance gap.

Our audit is limited to a single institution, anatomy, and architecture, and the
high-performance ceiling of cervical vertebral bodies and discs leaves little room
for disparity to emerge. Whether these findings generalise to harder anatomies or
lower-performing models remains open. We hope this work opens a rigorous
conversation about fairness in spine MRI segmentation.
